\documentclass[12pt]{article}

% Idioma y codificación
\usepackage[spanish, es-tabla]{babel}       %es-tabla para que se titule "Tabla"
\usepackage[utf8]{inputenc}

% Márgenes
\usepackage[a4paper,top=3cm,bottom=2.5cm,left=3cm,right=3cm]{geometry}

% Comentarios de bloque
\usepackage{verbatim}

% Paquetes de links
\usepackage[hidelinks]{hyperref}    % Permite enlaces
\usepackage{url}                    % redirecciona a la web

% Más opciones para enumeraciones
\usepackage{enumitem}


%------------------------------------------------------------------------

%Paquetes de figuras
\usepackage{caption}
\usepackage{subcaption} % Figuras al lado de otras
\usepackage{float}      % Poner figuras en el sitio indicado H.


% Paquetes de imágenes
\usepackage{graphicx}       % Paquete para añadir imágenes
\usepackage{transparent}    % Para manejar la opacidad de las figuras

% Paquete para usar colores
\usepackage[dvipsnames]{xcolor}
\usepackage{pagecolor}      % Para cambiar el color de la página

% Habilita tamaños de fuente mayores
\usepackage{fix-cm}


%------------------------------------------------------------------------

% Paquetes de matemáticas
\usepackage{amsmath, amsfonts, amssymb, mathrsfs}
\usepackage{cancel}     % Simplificar tachando
\usepackage{polynom}    % Divisiones y Ruffini

\usepackage{pgfplots}   %Representar funciones
\pgfplotsset{compat=1.18}  % Versión 1.18

%Definición de teoremas, etc.
\usepackage{amsthm}
%\swapnumbers   % Intercambia la posición del texto y de la numeración

\theoremstyle{plain}
\newtheorem{teo}{Teorema}[chapter]     % Se resetea en cada chapter
\newtheorem{coro}{Corolario}[teo]           % Se resetea en cada teorema
\newtheorem{prop}[teo]{Proposición}         % Usa el mismo contador que teorema
\newtheorem{lema}[teo]{Lema}                % Usa el mismo contador que teorema

\theoremstyle{remark}
\newtheorem*{observacion}{Observación}

\theoremstyle{definition}
\newtheorem{definicion}{Definición} [chapter]     % Se resetea en cada chapter
\newtheorem*{notacion}{Notacion}
\newtheorem*{ejemplo}{Ejemplo}
\newtheorem*{ejercicio*}{Ejercicio}             % No numerado
\newtheorem{ejercicio}{Ejercicio} [section]     % Se resetea en cada section


% Modificar el formato de la numeración del teorema "ejercicio"
\renewcommand{\theejercicio}{%
  \ifnum\value{section}=0 % Si no se ha iniciado ninguna sección
    \arabic{ejercicio}% Solo mostrar el número de ejercicio
  \else
    \thesection.\arabic{ejercicio}% Mostrar número de sección y número de ejercicio
  \fi
}


% \renewcommand\qedsymbol{$\blacksquare$}         % Cambiar símbolo QED
%------------------------------------------------------------------------

% Paquetes para encabezados
\usepackage{fancyhdr}
\pagestyle{fancy}
\fancyhf{}

\newcommand{\helv}{ % Modificación tamaño de letra
\fontfamily{}\fontsize{12}{12}\selectfont}
\setlength{\headheight}{15pt} % Amplía el tamaño del índice


%\usepackage{lastpage}   % Referenciar última pag   \pageref{LastPage}
\fancyfoot[C]{\thepage}

%------------------------------------------------------------------------

\begin{comment} % Descomentar en book. Comentar en article

% Conseguir que no ponga "Capítulo 1". Sino solo "1."
\makeatletter
\renewcommand{\@makechapterhead}[1]{%
  \vspace*{50\p@}%
  {\parindent \z@ \raggedright \normalfont
    \ifnum \c@secnumdepth >\m@ne
      \huge\bfseries \thechapter.\hspace{1em}\ignorespaces
    \fi
    \interlinepenalty\@M
    \Huge \bfseries #1\par\nobreak
    \vskip 40\p@
  }}
\makeatother
\end{comment}

%------------------------------------------------------------------------
% Personalizar la portada
\usepackage{titling}


%------------------------------------------------------------------------

% Comandos definidos
\newcommand{\bb}[1]{\mathbb{#1}}
\newcommand{\cc}[1]{\mathcal{#1}}
\newcommand{\Hop}{\stackrel{L'H\hat{o}pital}{=}}

\newcommand{\red}[1]{{\color{red}#1}} % Para integrales, destacar los cambios.

% Método de integración
\newcommand{\MetInt}[2]{
    \left[\begin{array}{c}
        #1 \\ #2
    \end{array}\right]
}

%------------------------------------------------------------------------


\begin{document}

    % 1. Foto de fondo
    % 2. Título
    % 3. Encabezado Izquierdo
    % 4. Color de fondo
    % 5. Coord x del titulo
    % 6. Coord y del titulo
    % 7. Fecha

    
    % 1. Foto de fondo
% 2. Título
% 3. Encabezado Izquierdo
% 4. Color de fondo
% 5. Coord x del titulo
% 6. Coord y del titulo
% 7. Fecha

\newcommand{\portada}[7]{

    \portadaBase{#1}{#2}{#3}{#4}{#5}{#6}{#7}
    \portadaBook{#1}{#2}{#3}{#4}{#5}{#6}{#7}
}

\newcommand{\portadaExamen}[7]{

    \portadaBase{#1}{#2}{#3}{#4}{#5}{#6}{#7}
    \portadaArticle{#1}{#2}{#3}{#4}{#5}{#6}{#7}
}




\newcommand{\portadaBase}[7]{

    % Tiene la portada principal y la licencia Creative Commons
    
    % 1. Foto de fondo
    % 2. Título
    % 3. Encabezado Izquierdo
    % 4. Color de fondo
    % 5. Coord x del titulo
    % 6. Coord y del titulo
    % 7. Fecha
    
    
    \thispagestyle{empty}               % Sin encabezado ni pie de página
    \newgeometry{margin=0cm}        % Márgenes nulos para la primera página
    
    
    % Encabezado
    \fancyhead[L]{\helv #3}
    \fancyhead[R]{\helv \nouppercase{\leftmark}}
    
    
    \pagecolor{#4}        % Color de fondo para la portada
    
    \begin{figure}[p]
        \centering
        \transparent{0.3}           % Opacidad del 30% para la imagen
        
        \includegraphics[width=\paperwidth, keepaspectratio]{assets/#1}
    
        \begin{tikzpicture}[remember picture, overlay]
            \node[anchor=north west, text=white, opacity=1, font=\fontsize{60}{90}\selectfont\bfseries\sffamily, align=left] at (#5, #6) {#2};
            
            \node[anchor=south east, text=white, opacity=1, font=\fontsize{12}{18}\selectfont\sffamily, align=right] at (9.7, 3) {\textbf{\href{https://losdeldgiim.github.io/}{Los Del DGIIM}}};
            
            \node[anchor=south east, text=white, opacity=1, font=\fontsize{12}{15}\selectfont\sffamily, align=right] at (9.7, 1.8) {Doble Grado en Ingeniería Informática y Matemáticas\\Universidad de Granada};
        \end{tikzpicture}
    \end{figure}
    
    
    \restoregeometry        % Restaurar márgenes normales para las páginas subsiguientes
    \pagecolor{white}       % Restaurar el color de página
    
    
    \newpage
    \thispagestyle{empty}               % Sin encabezado ni pie de página
    \begin{tikzpicture}[remember picture, overlay]
        \node[anchor=south west, inner sep=3cm] at (current page.south west) {
            \begin{minipage}{0.5\paperwidth}
                \href{https://creativecommons.org/licenses/by-nc-nd/4.0/}{
                    \includegraphics[height=2cm]{assets/Licencia.png}
                }\vspace{1cm}\\
                Esta obra está bajo una
                \href{https://creativecommons.org/licenses/by-nc-nd/4.0/}{
                    Licencia Creative Commons Atribución-NoComercial-SinDerivadas 4.0 Internacional (CC BY-NC-ND 4.0).
                }\\
    
                Eres libre de compartir y redistribuir el contenido de esta obra en cualquier medio o formato, siempre y cuando des el crédito adecuado a los autores originales y no persigas fines comerciales. 
            \end{minipage}
        };
    \end{tikzpicture}
    
    
    
    % 1. Foto de fondo
    % 2. Título
    % 3. Encabezado Izquierdo
    % 4. Color de fondo
    % 5. Coord x del titulo
    % 6. Coord y del titulo
    % 7. Fecha


}


\newcommand{\portadaBook}[7]{

    % 1. Foto de fondo
    % 2. Título
    % 3. Encabezado Izquierdo
    % 4. Color de fondo
    % 5. Coord x del titulo
    % 6. Coord y del titulo
    % 7. Fecha

    % Personaliza el formato del título
    \pretitle{\begin{center}\bfseries\fontsize{42}{56}\selectfont}
    \posttitle{\par\end{center}\vspace{2em}}
    
    % Personaliza el formato del autor
    \preauthor{\begin{center}\Large}
    \postauthor{\par\end{center}\vfill}
    
    % Personaliza el formato de la fecha
    \predate{\begin{center}\huge}
    \postdate{\par\end{center}\vspace{2em}}
    
    \title{#2}
    \author{\href{https://losdeldgiim.github.io/}{Los Del DGIIM}}
    \date{Granada, #7}
    \maketitle
    
    \tableofcontents
}




\newcommand{\portadaArticle}[7]{

    % 1. Foto de fondo
    % 2. Título
    % 3. Encabezado Izquierdo
    % 4. Color de fondo
    % 5. Coord x del titulo
    % 6. Coord y del titulo
    % 7. Fecha

    % Personaliza el formato del título
    \pretitle{\begin{center}\bfseries\fontsize{42}{56}\selectfont}
    \posttitle{\par\end{center}\vspace{2em}}
    
    % Personaliza el formato del autor
    \preauthor{\begin{center}\Large}
    \postauthor{\par\end{center}\vspace{3em}}
    
    % Personaliza el formato de la fecha
    \predate{\begin{center}\huge}
    \postdate{\par\end{center}\vspace{5em}}
    
    \title{#2}
    \author{\href{https://losdeldgiim.github.io/}{Los Del DGIIM}}
    \date{Granada, #7}
    \thispagestyle{empty}               % Sin encabezado ni pie de página
    \maketitle
    \vfill
}
    \portadaExamen{ffccA4.jpg}{Geometría II\\Examen II}{Geometría II. Examen II}{MidnightBlue}{-8}{28}{2023}{Arturo Olivares Martos}

    \begin{description}
        \item[Asignatura] Geometría II.
        \item[Curso Académico] 2022-23.
        \item[Grado] Doble Grado en Ingeniería Informática y Matemáticas.
        \item[Grupo] Único.
        \item[Profesor] Antonio Ros Mulero.
        \item[Descripción] Parcial del Tema 1 de Incidencias. Diagonalización.
        \item[Fecha] 12 de abril de 2023.
        \item[Duración] 60 minutos.
    
    \end{description}
    \newpage
    
    \begin{ejercicio} [6 ptos] Sea $f$ el endomorfismo de $\bb{R}^3$ cuya matriz asociada respecto de la base usual es
    \begin{equation*}
        A = \left( \begin{array}{ccc}
            0 & 0 & 1 \\
            -a & 1-a & -a\\
            a & 0 & a-1
        \end{array}\right) \qquad  a\in \bb{R}.
    \end{equation*}
    \begin{enumerate}
        \item Encontrar los valores de $a$ para los que $f$ es diagonalizable.


        Calculamos el polinomio característico de $A$:
        \begin{equation*}\begin{split}
                P_{A}(\lambda) & = |A-\lambda_0 I| = \left| \begin{array}{ccc}
                -\lambda_0 & 0 & 1 \\
                -a & 1-a-\lambda_0 & -a \\
                a & 0 & a-1-\lambda_0 \\
                \end{array}\right|
                =
                (1-a-\lambda_0) \left| \begin{array}{cc}
                -\lambda_0 & 1 \\
                a & a-1-\lambda_0 \\
                \end{array}\right|
                =\\&
                =
                (1-a-\lambda_0) \left| \begin{array}{cc}
                -\lambda_0 & 1+\lambda_0 \\
                a & -1-\lambda_0 \\
                \end{array}\right|
                =
                (1-a-\lambda_0)(1+\lambda_0)\left| \begin{array}{cc}
                -\lambda_0 & 1 \\
                a & -1 \\
                \end{array}\right|
                =\\&=
                (1-a-\lambda_0)(1+\lambda_0)(\lambda_0-a)
            \end{split}\end{equation*}

            Por tanto, los valores propios son: $\{1-a, -1, a\}$:
            \begin{equation*}
                1-a=a\Longleftrightarrow a=\frac{1}{2} 
                \qquad 1-a=-1\Longleftrightarrow a=2
                \qquad -1=a
            \end{equation*}

            Por tanto, distinguimos según los valores de $a$:
            \begin{itemize}
                \item \underline{Para $a\neq -1, \frac{1}{2}, 2$}:

                Tenemos que tiene tres valores propios distintos, por lo que es diagonalizable.

                \item \underline{Para $a=-1$}:

                Tenemos que la multiplicidad algebraica de $\lambda = -1$ es doble. Veamos su multiplicidad geométrica:
                \begin{equation*}
                    V_{-1} = \left\{x\in\bb{R}^3 \left| \begin{array}{c}
                         x_1+x_3=0  \\
                         x_1+3x_2+x_3=0 \\
                         -x_1-x_3=0
                    \end{array}\right.\right\} = \mathcal{L}\left(\left\{ \left(\begin{array}{c}
                            1 \\
                            0 \\
                            -1 \\
                       \end{array}\right)
                       \right\}\right)
                \end{equation*}

                Por tanto, como la multiplicidad geométrica es distinta a la algebraica $(1\neq 2)$, tenemos que no es diagonalizable.

                \item \underline{Para $a=\frac{1}{2}$}:

                Tenemos que la multiplicidad algebraica de $\lambda = \frac{1}{2}$ es doble. Veamos su multiplicidad geométrica:
                \begin{equation*}
                    V_{\frac{1}{2}} = \left\{x\in\bb{R}^3 \left| \begin{array}{c}
                         -\frac{1}{2}x_1+x_3=0  \\
                         -\frac{1}{2}x_1-\frac{1}{2}x_3=0 \\
                         \frac{1}{2}x_1-x_3=0
                    \end{array}\right.\right\} = \mathcal{L}\left(\left\{ \left(\begin{array}{c}
                            0 \\
                            1 \\
                            0 \\
                       \end{array}\right)
                       \right\}\right)
                \end{equation*}

                Por tanto, como la multiplicidad geométrica es distinta a la algebraica $(1\neq 2)$, tenemos que no es diagonalizable.


                \item \underline{Para $a=2$}:

                Tenemos que la multiplicidad algebraica de $\lambda = -1$ es doble. Veamos su multiplicidad geométrica:
                \begin{equation*}
                    V_{-1} = \left\{x\in\bb{R}^3 \left| \begin{array}{c}
                         x_1+x_3=0  \\
                         -2x_1-2x_3=0 \\
                         2x_1+2x_3=0
                    \end{array}\right.\right\} = \mathcal{L}\left(\left\{ \left(\begin{array}{c}
                            0 \\
                            1 \\
                            0 \\
                       \end{array}\right),
                       \left(\begin{array}{c}
                            1 \\
                            0 \\
                            -1 \\
                       \end{array}\right)
                       \right\}\right)
                \end{equation*}

                Por tanto, como la multiplicidad geométrica es igual a la algebraica $(2)$, tenemos que sí es diagonalizable.
            \end{itemize}

            Por tanto, tenemos que $A$ es diagonalizable $f\forall a\neq \{-1, \frac{1}{2}\}$.

            
        \item Para $a=1$, estudiar si las matrices $A$ y $A^2$ son semejantes y diagonalizar alguna de ellas, si ello es posible.

        Diagonalizamos en primer lugar $A$. Tenemos que tiene tres valores propios distintos $\{-1, 0, 1\}$, por lo que calculamos los subespacios propios:
        \begin{equation*}
            V_{-1} = \left\{x\in\bb{R}^3 \left| \begin{array}{c}
                 x_1+x_3=0  \\
                 -x_1+x_2-x_3=0 \\
                 x_1+x_3=0
            \end{array}\right.\right\} = \mathcal{L}\left(\left\{ \left(\begin{array}{c}
                    1 \\
                    0 \\
                    -1 \\
               \end{array}\right)
               \right\}\right)
        \end{equation*}
        \begin{equation*}
            V_{0} = \left\{x\in\bb{R}^3 \left| \begin{array}{c}
                 x_3=0  \\
                 -x_1-x_3=0 \\
                 x_1=0
            \end{array}\right.\right\} = \mathcal{L}\left(\left\{ \left(\begin{array}{c}
                    0 \\
                    1 \\
                    0 \\
               \end{array}\right)
               \right\}\right)
        \end{equation*}
        \begin{equation*}
            V_{1} = \left\{x\in\bb{R}^3 \left| \begin{array}{c}
                 -x_1+x_3=0  \\
                 -x_1-x_2-x_3=0 \\
                 x_1-x_3=0
            \end{array}\right.\right\} = \mathcal{L}\left(\left\{ \left(\begin{array}{c}
                    -1 \\
                    2 \\
                    -1 \\
               \end{array}\right)
               \right\}\right)
        \end{equation*}

        Por tanto, tenemos que $D=P^{-1}AP$, con:
        \begin{equation*}
            P = \left( \begin{array}{ccc}
                0 & -1 & 1 \\
                1 & 2 & 0\\
                0 & -1 & -1
            \end{array}\right)
            \hspace{2cm}
            D = \left( \begin{array}{ccc}
                0 \\
                &1\\
                &&-1
            \end{array}\right)
        \end{equation*}

        Por tanto, $A=PDP^{-1}$. Calculamos ahora $A^2$:
        \begin{equation*}
            A^2 = PDP^{-1}\cdot PDP^{-1} = PD^2P^{-1}
        \end{equation*}

        Tenemos que:
        \begin{equation*}
            A\sim D=\left( \begin{array}{ccc}
                0 \\
                &1\\
                &&-1
            \end{array}\right)
            \hspace{2cm}
            A^2\sim D^2=\left( \begin{array}{ccc}
                0 \\
                &1\\
                &&1
            \end{array}\right)
        \end{equation*}

        Como tanto, como $D\neq D^2$, tenemos que $A\nsim A^2$.
    \end{enumerate}
\end{ejercicio}

\vspace{1cm}

\begin{ejercicio} [4 ptos]
    Razonar si las siguientes afirmaciones son verdaderas o falsas:
    \begin{enumerate}
        \item Sobre el cuerpo de los complejos, si la matriz cuadrada $A$ de orden 4 verifica $|A|=-1$, entonces $A$ es diagonalizable.

        \begin{equation*}
            A=\left( \begin{array}{cccc}
                1 & 1 & 0 & 0 \\
                0 & 1 & 0 & 0 \\
                0 & 0 & -1 & 0 \\
                0 & 0 & 0 & 1 \\
            \end{array}\right)
        \end{equation*}

        Tenemos que $|A|=-1$, y calculemos su polinomio característico:
        \begin{equation*}\begin{split}
            P_{A}(\lambda) & = |A-\lambda_0 I| = \left| \begin{array}{cccc}
                1-\lambda_0 & 1 & 0 & 0 \\
                0 & 1-\lambda_0 & 0 & 0 \\
                0 & 0 & -1-\lambda_0 & 0 \\
                0 & 0 & 0 & 1-\lambda_0 \\
            \end{array}\right|
            =
            -(1-\lambda_0)^3(1+\lambda_0)
        \end{split}\end{equation*}

        Por tanto, tenemos que los valores propios son $\lambda_0=\{-1, 1\}$. Calculamos la multiplicidad geométrica de $\lambda_0=1$.
        \begin{equation*}
            V_{1} = \left\{x\in\bb{R}^4 \left| \begin{array}{c}
                 x_2=0  \\
                 -2x_3=0 \\
            \end{array}\right.\right\}
        \end{equation*}

        Por tanto, tenemos que la multiplicidad algebraica del $1$ es $3$, pero su multiplicidad geométrica es $2$. Por tanto, tenemos que no es diagonalizable.

        El enunciado, por tanto, es falso.
        
        

        \item Si una matriz $A\in \mathcal{M}_n(\bb{R})$ verifica $A^3+A=0$, entonces $A$ es diagonalizable.

        Tenemos que $A^3=-A$. Sea el contraejemplo siguiente:
        \begin{equation*}
            A=\frac{1}{2}\left(\begin{array}{cc}
                1 & -\sqrt{3} \\
                \sqrt{3} & 1
            \end{array}\right)
        \end{equation*}

        Tenemos que no es digonalizable pero se cumple que $A^3=-Id$.

        \begin{observacion}
            Se trata de un giro de $\theta = \frac{\pi}{3}$ en el plano.
        \end{observacion}
    \end{enumerate}
\end{ejercicio}



\end{document}